\subsection{A concrete descent obstruction in the nonstandard cubic algebra}

The general nonstandard setting is not an ordinary Laurent-polynomial
version of the standard canonical-basis problem. The following obstruction
uses the actual cubic example in GCT VII--VIII, rather than an unrelated
positive-semidefinite countermodel.

\paragraph{Exact known example.}
In GCT VII, Section~7.1.1--7.1.2, printed pp.~18--20,
equations~(19)--(24), let $B_0$ be the nonstandard algebra for
$H=GL_2$, $X=\operatorname{Sym}^3\mathbb C^2$, and tensor degree three,
over $F=\mathbb C(q)$. It has dimension $21$ and a basis of the
alternating monomials $Q_\sigma$ of length at most ten, including
the empty monomial. The displayed splitting extension is
\[
 K=F(x),\qquad x^2=D(q),\qquad
 D(q)=(q^8+1)^2q^{24}A(q),
 \tag{P-F36a}
\]
where
\[
 A(q)=5q^{16}+8q^{12}-4q^{10}+18q^8-4q^6+8q^4+5.
 \tag{P-F36b}
\]
Over $K$, the algebra is a direct sum of one one-dimensional
simple block and five four-dimensional simple blocks.
The blocks $U_1,U_2$ are interchanged by $x\mapsto-x$;
see equation~(22), the definitions following it, and equation~(23).
These are explicit assertions for the worked example, not an
assumption that the conjectural algorithms work in every degree.

GCT VIII repeats the discriminant in Section~5.3, equation~(18),
printed p.~31. In its Figure~6, p.~38, the coefficient of
$Q_{1212121}$ in the displayed element $u_1^1$ is
\[
 c(q)=-\frac{P(q)+x}{2q^{20}},
 \tag{P-F36c}
\]
where
\[
 \begin{split}
 P(q)={}&2q^{30}+3q^{28}+2q^{26}+6q^{24}-2q^{22}
          +22q^{20}-2q^{18}\\
       &+6q^{16}+2q^{14}+3q^{12}+2q^{10}.
 \end{split}
\]
The corresponding element $u_2^1$ is obtained by replacing $x$ by
$-x$. This coefficient is recorded in the monomial basis just
specified; it is not asserted here to be a coefficient in a
different ambient canonical basis.

\begin{proposition}[P-F36: power regrading and rational projectors do not remove the quadratic datum]
\label{prop:nonstandarddescent}
For the displayed nonstandard cubic algebra:
\begin{enumerate}
\item For every positive integer $k$, the rational function $D(t^k)$
is not a square in $\mathbb C(t)$. Consequently, after the power
regrading $q=t^k$, the coefficient $c(t^k)$ is not in
$\mathbb C(t)$. In particular $c(q)$ is not the character of any
finite rationally graded vector space, even allowing arbitrary
complex coefficients or virtual signs. After any power regrading,
no nonzero rational-function scalar normalization makes the
coefficient a finite Laurent character.
\item There is no $F$-linear projector on $B_0$ whose scalar
extension has image exactly $U_1$ or exactly $U_2$. More generally,
if a graded $B_0$-bimodule or the cohomology of a complex of
$B_0$-bimodules is defined over $F$, then after extension to $K$
the multiplicities of $U_1$ and $U_2$ as simple bimodules are equal
in each degree. Thus neither block alone is a descended graded
object of that type.
\end{enumerate}
The same descent conclusions hold after any power regrading.
\end{proposition}

\begin{proof}
We first prove that $A(q)$ is not a square even in $\mathbb C(q)$.
If the square of a rational function is a polynomial, that rational
function has no finite pole and is itself a polynomial. Suppose
therefore that $A(q)=R(q)^2$ with $R\in\mathbb C[q]$.
Since $A$ is even, the integral-domain identity
\[
 (R(-q)-R(q))(R(-q)+R(q))=0
\]
makes $R$ either even or odd. Its constant term is nonzero because
$A(0)=5$, so it is even. Its degree is eight. Write
\[
 R(q)=a q^8+b q^6+c_0 q^4+d q^2+e.
\]
The coefficients at degrees $16,14,12,10$ respectively give
\[
 a^2=5,\qquad b=0,\qquad c_0=4/a,\qquad d=-2/a.
\]
The coefficient at degree eight then gives
\[
 2ae+c_0^2=18,
 \qquad ae=\frac{37}{5}.
\]
But the constant term gives $e^2=5$, and hence
$(ae)^2=a^2e^2=25$, contrary to $ae=37/5$.
Thus $A$ is not a square.

Over $\mathbb C$, a polynomial is a square up to a nonzero scalar
exactly when all its roots have even multiplicity. The scalar
itself has a square root. Since $A$ is not a square and $A(0)\ne0$,
there is a nonzero root $\alpha$ of odd multiplicity. For every
$k\ge1$, each root of $t^k=\alpha$ is nonzero and the derivative
$kt^{k-1}$ is nonzero there. The corresponding root of $A(t^k)$
therefore has the same odd multiplicity. It follows that
$A(t^k)$ is not a square in $\mathbb C(t)$. Multiplication by the
explicit nonzero square factor in (P-F36a) does not change its
square class, so $D(t^k)$ is not a square either.

In the quadratic extension obtained by adjoining a square root
$x_k$ of $D(t^k)$, formula~(P-F36c) reads
\[
 c(t^k)=-\frac{P(t^k)+x_k}{2t^{20k}}.
\]
Its rational coefficient of $x_k$ is nonzero. If this expression
belonged to $\mathbb C(t)$, solving for $x_k$ would put the
square root in that field, a contradiction. Multiplying it by
a nonzero element of $\mathbb C(t)$ cannot change this conclusion.

A finite rationally graded vector space has a character that
becomes a finite Laurent polynomial in $t$ after clearing the
denominators of its grading. A virtual difference gives a signed
Laurent polynomial; allowing arbitrary scalar coefficients still
gives a Laurent polynomial. All lie in $\mathbb C(t)$.
This proves the first assertion. It concerns finite Laurent
characters after power regrading, not arbitrary local Puiseux
series: algebraic functions can have such local expansions.

For the second assertion, let $\tau$ be the nontrivial automorphism
of $K/F$. Scalar extension of any $F$-linear map commutes with
the natural semilinear action of $\tau$. Its image is therefore
$\tau$-stable. But $\tau(U_1)=U_2$ and the two blocks are distinct
summands in the displayed direct decomposition. Hence that image
cannot be either block alone. This already proves the projector
claim without imposing bimodule linearity on the projector.

For the more general claim, scalar extension of an $F$-defined
$B_0$-bimodule carries this semilinear action compatibly with
the algebra. The split semisimple algebra $B_0\otimes_F K$
has distinct simple bimodules $U_i$. Conjugating a decomposition
interchanges the $U_1$ and $U_2$ multiplicities. Since the module
itself is unchanged up to the natural semilinear identification,
uniqueness of the simple decomposition makes these multiplicities
equal. Apply the same reasoning separately in each preserved
grading degree.

Extension by a field is exact. For any complex of $F$-defined
bimodules its cohomology after extension is the extension of its
cohomology before extension. Thus the same argument applies to
every cohomological degree, and also to signed Euler classes.
Finally, the first part proves that the quadratic extension
remains nontrivial after every power regrading, so the entire
descent argument repeats over $\mathbb C(t)$.
\end{proof}

\paragraph{The sharper transfer requirement.}
A direct purity or determinant construction whose output is an
ordinary finite graded multiplicity character cannot identify
the displayed coefficient, even after a fractional grading or
rational scalar normalization. Likewise a construction that stays
over the rational function field and yields an actual bimodule
projector cannot select the two conjugate blocks separately.
It must retain the quadratic extension and its descent datum,
or work with a conjugate sum and then provide an additional
comparison that recovers the desired separate objects.

This does not refute arbitrary theta-lift constructions or
algebraic coefficient systems: those can carry nontrivial field
extensions. It also does not refute GCT VIII's relaxed conjectures,
which explicitly concern coefficients of minimal polynomials,
with permitted signs and factors, rather than identifying every
raw algebraic coefficient with a positive Laurent polynomial;
see Conjectures~2.8--2.9, pp.~16--17, and 3.5--3.6, pp.~24--25.
The supplied Alp\"oge \emph{Remark}, pp.~3--4, steps~4, 8--12,
does not specify a comparison of its automorphic projectors with
these two nonstandard blocks. P-F36 states a necessary condition
on such a comparison; it makes no assertion that the cited
theta lifts lack the required extension or local systems.

The exact primary sources are
\url{https://arxiv.org/pdf/0709.0749}, Section~7.1.1--7.1.2,
pp.~18--20, equations~(19)--(24), and
\url{https://arxiv.org/pdf/0709.0751}, Section~5.3, p.~31,
equation~(18), and Figure~6, p.~38.
